% !TeX root = prueba02_version01.tex
% !TeX program = pdflatex
% Prueba 2: versión 1, dificultad media, 120 minutos, 10 puntos.
\documentclass[10pt,a5paper]{article}
\newif\ifmostrarSoluciones
\mostrarSolucionestrue
\usepackage{estilo_apuntes}
\dosPorCarillafalse % Use \dosPorCarillatrue para imprimir dos copias por carilla.
\encabezadoprueba{Prueba 2}{Versión 1}

\begin{document}
\begin{center}
  {\color{Azul}\large\bfseries Prueba 2 · Versión 1}
\end{center}

\noindent\textbf{Nombre:} \hrulefill

\medskip
\noindent\textbf{Indicaciones}
\begin{itemize}
  \item La duración de la prueba es de 2 horas. Se califica sobre 10 puntos.
  \item Escriba el desarrollo de manera clara y ordenada. Justifique las respuestas e indique las propiedades y restricciones utilizadas. Al aplicar un teorema, verifique sus hipótesis y exprese su conclusión. Una respuesta sin justificación recibirá únicamente crédito parcial.
  \item Resuelva con las herramientas de límites y continuidad estudiadas, sin derivadas ni regla de L'Hôpital. No se solicitan demostraciones de teoremas.
  \item No se permite el uso de dispositivos electrónicos. Cualquier intento de copia penalizará la evaluación con una calificación total de cero.
\end{itemize}

\noindent\textbf{Ejercicios}
\begin{enumerate}
  \item \textbf{(4 puntos)} Considere
  \[
    R(x)=\frac{x^2-1}{(x-1)(x-3)},\qquad
    D_R=\mathbb{R}\setminus\{1,3\},
  \]
  y
  \[
    H(x)=\frac{\sqrt{9x^2+4}}{x+1},\qquad
    D_H=\mathbb{R}\setminus\{-1\}.
  \]
  \begin{enumerate}
    \item \textbf{(2 puntos)} Calcule el límite de $R$ en $1$ y sus límites laterales en $3$. Determine todas las asíntotas verticales de $R$.
    \item \textbf{(2 puntos)} Calcule los límites de $H$ cuando $x\to+\infty$ y $x\to-\infty$. Determine sus asíntotas horizontales e indique a qué extremo corresponde cada una.
  \end{enumerate}

  \item \textbf{(3 puntos)} Para $a,b\in\mathbb{R}$, sea
  \[
    F_{a,b}(x)=\begin{cases}
      \dfrac{x^2-4}{x-2},&x<2,\\[2mm]
      ax+b,&2\leq x\leq4,\\
      \sqrt{x+5},&x>4.
    \end{cases}
  \]
  \begin{enumerate}
    \item \textbf{(2 puntos)} Determine $a,b$ para que $F_{a,b}$ sea continua en $\mathbb{R}$. Verifique los puntos de unión y explique qué ocurre en el interior de cada tramo.
    \item \textbf{(1 punto)} Ahora tome $a=0$, $b=4$. Clasifique la discontinuidad en $x=4$ e indique si se puede eliminar cambiando únicamente $F_{0,4}(4)$.
  \end{enumerate}

  % Adaptación de Bartle--Sherbert, Introduction to Real Analysis, 4.ª ed., §5.3, ej. 5.
  % Se omite la aproximación con calculadora y se incorpora máximo--mínimo.
  \item \textbf{(3 puntos)} Considere $p(x)=x^4+7x^3-9$ en $[-8,2]$.
  \begin{enumerate}
    \item \textbf{(1 punto)} Justifique que $p$ alcanza un máximo y un mínimo absolutos en $[-8,2]$, indicando el teorema y verificando sus hipótesis.
    \item \textbf{(1,5 puntos)} Justifique que $p$ tiene al menos dos ceros distintos: uno en $(-8,-7)$ y otro en $(1,2)$. No se pide calcularlos ni aproximarlos.
    \item \textbf{(0,5 puntos)} Sin hallar el punto de mínimo, justifique que el valor mínimo es negativo y que se alcanza en algún punto de $(-8,2)$.
  \end{enumerate}
\end{enumerate}

\ifmostrarSoluciones
\newpage
\begin{solucion}[de la versión 1]
\begin{enumerate}
  \item
  \begin{enumerate}
    \item Para $x\neq1,3$, $R(x)=(x+1)/(x-3)$. Por tanto,
    \[
      \lim_{x\to1}R(x)=-1.
    \]
    Cerca de $3$ el numerador es positivo y tiende a $4$. El denominador es negativo a la izquierda y positivo a la derecha. Así,
    \[
      \lim_{x\to3^-}R(x)=-\infty,\qquad
      \lim_{x\to3^+}R(x)=+\infty.
    \]
    La única asíntota vertical es $x=3$. El punto $1$ permanece excluido del dominio, pero el límite allí es finito. En los demás puntos, el denominador no se anula.
    \item Para $x\neq0$,
    \[
      H(x)=\frac{|x|}{x}\,
      \frac{\sqrt{9+4/x^2}}{1+1/x}.
    \]
    El factor $|x|/x$ vale $1$ para $x>0$ y $-1$ para $x<0$. Por ello,
    \[
      \lim_{x\to+\infty}H(x)=3,\qquad
      \lim_{x\to-\infty}H(x)=-3.
    \]
    Las asíntotas horizontales son $y=3$ hacia $+\infty$ y $y=-3$ hacia $-\infty$.
  \end{enumerate}

  \item
  \begin{enumerate}
    \item El primer tramo coincide con $x+2$ para $x<2$. La continuidad en $2$ exige
    \[
      4=2a+b=F_{a,b}(2).
    \]
    La continuidad en $4$ exige
    \[
      F_{a,b}(4)=4a+b=\sqrt{4+5}=3.
    \]
    Resolviendo simultáneamente, $2a=-1$, de donde
    \[
      a=-\frac12,\qquad b=5.
    \]
    En los interiores de los tramos hay funciones polinómicas o una raíz de radicando positivo, por lo que son continuas allí. Las dos igualdades verifican la continuidad en los puntos de unión y, por tanto, en todo $\mathbb{R}$.
    \item Para $a=0$, $b=4$,
    \[
      \lim_{x\to4^-}F_{0,4}(x)=4,\qquad
      \lim_{x\to4^+}F_{0,4}(x)=3.
    \]
    Es una discontinuidad de salto. Cambiar solamente el valor en $4$ no modifica los límites laterales y no elimina el salto.
  \end{enumerate}

  \item
  \begin{enumerate}
    \item La función $p$ es polinómica, por lo que es continua en todo $[-8,2]$. Como este intervalo es cerrado y acotado, el teorema de máximo--mínimo garantiza que ambos extremos absolutos se alcanzan.
    \item Se calculan los valores
    \[
      p(-8)=503>0,\qquad p(-7)=-9<0,
    \]
    \[
      p(1)=-1<0,\qquad p(2)=63>0.
    \]
    La función es continua en $[-8,-7]$ y en $[1,2]$. Bolzano garantiza un cero en $(-8,-7)$ y otro en $(1,2)$. Son distintos porque los intervalos son disjuntos. Esto establece al menos dos ceros, sin afirmar que sean los únicos.
    \item Sea $m$ el mínimo absoluto, cuya existencia ya se justificó. Como $0\in[-8,2]$ y $p(0)=-9$, se tiene $m\leq-9<0$. Los extremos del intervalo tienen valores positivos, $p(-8)=503$ y $p(2)=63$, así que ninguno puede ser un punto de mínimo. Por tanto, el mínimo se alcanza en el interior $(-8,2)$.
  \end{enumerate}

\end{enumerate}
\end{solucion}
\fi

\end{document}
